% ---------------------------------------------------------------------
%  TÓM TẮT
% ---------------------------------------------------------------------
\begin{abstract}

\begin{chuy}
\goiy{tóm tắt khoảng 200 đến 300 từ, đi theo năm ý và mỗi ý một đến hai câu:
bối cảnh, vấn đề cần giải, cách em giải, kết quả chính kèm con số, và ý
nghĩa. Đây là phần hội đồng đọc đầu tiên, nên viết sau cùng khi đã có kết
quả thật. Không đưa trích dẫn và không đưa nội dung chưa có trong báo cáo.}
\end{chuy}

\begin{mau}
Nhận dạng loài vật qua ảnh là bài toán có nhiều ứng dụng trong giám sát đa
dạng sinh học và giáo dục. Thách thức nằm ở chỗ ảnh chụp ngoài thực địa có
nền phức tạp, điều kiện ánh sáng thay đổi và nhiều loài trông gần giống
nhau. Đồ án xây dựng một hệ thống nhận dạng loài vật dựa trên mạng nơ-ron
tích chập, kết hợp học chuyển giao từ mô hình đã tiền huấn luyện và các phép
tăng cường dữ liệu phù hợp với đặc thù ảnh thực địa. Hệ thống được huấn
luyện và đánh giá trên tập dữ liệu gồm 12.000 ảnh thuộc 30 loài, đạt độ
chính xác 91,2\% trên tập kiểm thử độc lập, cao hơn 6,4 điểm phần trăm so
với mô hình huấn luyện từ đầu. Toàn bộ mô hình được triển khai thành một ứng
dụng web cho phép người dùng tải ảnh lên và nhận kết quả kèm độ tin cậy
trong dưới một giây. Kết quả cho thấy học chuyển giao là lựa chọn hợp lý khi
dữ liệu huấn luyện ở quy mô vừa, và hệ thống đủ nhẹ để chạy trên máy tính
cá nhân thông thường.

\vspace{0.8em}
\noindent\textbf{Từ khoá:} mạng nơ-ron tích chập, học chuyển giao, nhận dạng
ảnh, thị giác máy tính.
\end{mau}

\end{abstract}
